\documentclass[11pt,a4paper]{article}
\usepackage[T1]{fontenc}
\usepackage{lmodern,amsmath,amssymb,amsthm,mathtools}
\usepackage[left=25mm,right=25mm,top=24mm,bottom=24mm,headheight=14pt]{geometry}
\usepackage{booktabs,array,microtype,fancyhdr,url}
\usepackage[dvipsnames]{xcolor}
\usepackage[colorlinks=true,linkcolor=MidnightBlue,citecolor=MidnightBlue,urlcolor=MidnightBlue]{hyperref}
\usepackage{bookmark}
\newcolumntype{P}[1]{>{\raggedright\arraybackslash}p{#1}}
\renewcommand{\labelitemi}{$\bullet$}
\expandafter\def\expandafter\UrlBreaks\expandafter{\UrlBreaks\do\_\do\a\do\b\do\c\do\d\do\e\do\f\do\g\do\h\do\i\do\j\do\k\do\l\do\m\do\n\do\o\do\p\do\q\do\r\do\s\do\t\do\u\do\v\do\w\do\x\do\y\do\z\do\A\do\B\do\C\do\D\do\E\do\F\do\G\do\H\do\I\do\J\do\K\do\L\do\M\do\N\do\O\do\P\do\Q\do\R\do\S\do\T\do\U\do\V\do\W\do\X\do\Y\do\Z}
\hypersetup{pdftitle={Certified Connection from Local Period-Doubling to the Captured Alternating Cycle in the Native Triad},pdfauthor={Hilmir Frímann Halldórsson (pzychozen)},pdfsubject={Research edition v0.1; project-internal manuscript review complete; author-authorized for repository publication; manuscript dated 12 November 2026},pdfkeywords={native triad, period doubling, validated continuation, alternating two-cycle, Paper G companion}}
\setlength{\parindent}{0pt}\setlength{\parskip}{5pt}
\setlength{\emergencystretch}{2em}
\raggedbottom\clubpenalty=10000\widowpenalty=10000\displaywidowpenalty=10000
\pagestyle{fancy}\fancyhf{}
\fancyhead[L]{\small Certified alternating-branch connection}
\fancyhead[R]{\small Research edition v0.1}\fancyfoot[C]{\thepage}
\renewcommand{\headrulewidth}{0.3pt}
\numberwithin{equation}{section}
\newtheorem{theorem}{Theorem}[section]
\newtheorem{lemma}[theorem]{Lemma}
\newtheorem{proposition}[theorem]{Proposition}
\newtheorem{corollary}[theorem]{Corollary}
\newcommand{\R}{\mathbb R}\newcommand{\C}{\mathbb C}
\newcommand{\Fh}{\widehat F}\newcommand{\eps}{\varepsilon}
\newcommand{\iu}{\mathrm i}\newcommand{\one}{\mathbf 1}
\newcommand{\diag}{\operatorname{diag}}\newcommand{\Argz}{\operatorname{Arg}_0}
\newcommand{\Span}{\operatorname{span}}\newcommand{\Id}{\mathrm I}
\newcommand{\entC}{\mathrm C}\newcommand{\entS}{\mathrm S}
\newcommand{\normr}[1]{\left\|#1\right\|_r}
\begin{document}
\thispagestyle{plain}
{\LARGE\bfseries Certified Connection from Local Period-Doubling to the Captured Alternating Cycle in the Native Triad\par}
\vspace{5pt}
{\small A companion to \emph{Geometry-Attached Axial Observables, Period-Doubling, and Certified Entrance Capture in the Tri-Octagon Map}\par}
\vspace{7pt}
{\large Hilmir Fr\'imann Halld\'orsson (pzychozen)\par}
{\small Research edition v0.1 \quad 12 November 2026\par}
\vspace{5pt}
{\small\textbf{Status.} Project-internal manuscript review complete; author-authorized for repository publication. No external peer review or proof-assistant verification is claimed.\par}

\begin{abstract}
Paper G established a local supercritical period-doubling branch and, separately, capture by an alternating cycle at a specified finite parameter. We explain the accepted connection between those particular objects for the unchanged exact native-triad family. An amplitude formulation removes the apparent singularity at onset and fixes a nontrivial phase orientation. Its first uniform patch includes the bifurcation parameter. A finite certificate then covers $0<\mu\le1/200$ with 128 patches and 127 common-uniqueness joins. The entire refined endpoint enclosure lies strictly inside the exact inherited M3 trap, identifying the continued root with its unique two-step fixed point. The branch is a graph over $\mu$, continuous to the host and smooth for positive parameters. Inherited symmetry lifts every member to a primitive full-state two-cycle. A separate derivative bound proves local attraction modulo common phase only on $4999/1000000\le\mu\le1/200$. Attraction throughout the connected interval remains unresolved. These are mathematical statements about a fixed discrete map, with no physical-field identification or parameter ramp imposed on an orbit.
\end{abstract}

\section{The missing connection}\label{sec:introduction}
Existence near a bifurcation and capture at a distant parameter answer different questions. Paper G's M2 theorem supplies a locally attracting alternating branch for sufficiently small positive $\mu$. Its repaired M3 theorem supplies a captured two-cycle at the exact value $\mu=1/200$. Neither theorem, by itself, identifies that captured cycle with the local branch. Paper G explicitly retained this non-continuation boundary \cite{PG}. The present companion supplements that correct earlier edition; it does not revise its equations or evidence.

The accepted investigation closes precisely this gap \cite{AB,Review}. Its proof must reach the actual onset, keep roots identified while the parameter varies, and identify the whole endpoint enclosure with the exact M3 trap. A small residual or a plausible amplitude plot would not provide those steps. The contribution is the connection of previously established objects, rather than discovery of period doubling or validated continuation methods. The retained prior-work comparison is a bounded comparison within the inspected project corpus, not a worldwide priority claim. The different turning-to-alignment family and its update-50 certificate concern initial-condition evolution toward an equilibrium, not this parameter branch \cite{TC}.

\begin{table}[htbp]\centering\small
\begin{tabular}{P{.13\linewidth}P{.48\linewidth}P{.29\linewidth}}\toprule
Record & Mathematical content & Status in this companion\\\midrule
I01 & M2 host, compression, local flip and corrected normalization & Inherited accepted premise\\
I02 & M3 trap, unique endpoint, H exchange and repaired capture & Inherited accepted premise\\
N01/N02 & Actual onset anchor, connected graph and endpoint identity & New accepted connection\\
N03 & Primitive full-state two-cycles and scoped sign reversal & Inherited symmetry applied to the connected family\\
N04 & Direct two-step attraction bound on the explicit endpoint strip & Separate narrow result\\\bottomrule
\end{tabular}
\caption{Contribution and premise boundary. All six source-ledger records are mapped to final manuscript locations in the supplied crosswalk.}\label{tab:contributions}
\end{table}

\section{Exact family and inherited premises}\label{sec:family}
\subsection{The simultaneous ordered map}\label{sec:map}
For $\Omega\in\C^3$, the native update first forms the complete prestage $v$, then applies the phase kicks:
\begin{align}
v_i&=\Omega_i+\eps\Omega_i(k_i-|\Omega_i|^2)
       +g\sum_{j\ne i}(\Omega_j-\Omega_i),\label{eq:native-pre}\\
\delta_i&=\lambda\sum_{j\ne i}
 \sin\!\left(3[\Argz v_j-\Argz v_i]\right),
 &F_\lambda(\Omega)_i&=v_i e^{\iu\delta_i}.\label{eq:native-phase}
\end{align}
Every prestage uses the same old state and every kick uses the same complete prestage. The convention is $\Argz(0)=0$. The implementation separately returns the prestage unchanged when $\lambda=0$; our family has positive $\lambda$ and never needs that special branch. No amplitude normalization, sequential channel update, extra state, receiver, feedback or physical time is added.

Throughout the connection, fix
\begin{equation}
\eps=\frac1{20},\qquad g=\frac15,\qquad
k=(1,1.2208964704604097,6.35310346037241),
\quad\lambda=\lambda_c+\mu.\label{eq:parameters}
\end{equation}
The terminating decimals in $k$ are exact rationals. Let $L=\one\one^T-3\Id_3$ and let $u$ be the positive root selected by the accepted M2 box:
\begin{equation}
\eps u\odot(k-u^2)+gLu=0.\label{eq:host}
\end{equation}
The host is this exact root, not its printed center or the midpoint of an enclosure. The inherited Krawczyk image encloses it and is an accepted numerical input. No new host theorem is required here.

Set $a_{12}=(u_A^2+u_B^2)^{1/2}$ and $r_u=\|u\|_2$. The source basis is
\begin{equation}
v_1=\frac{(u_B,-u_A,0)^T}{a_{12}},\qquad
v_2=\frac{(u_Au_C,u_Bu_C,-a_{12}^2)^T}{a_{12}r_u},
\qquad V=(v_1\ v_2).\label{eq:basis}
\end{equation}
Thus $V^TV=\Id_2$ and $V^Tu=0$. All occurrences of $u$, $V$ and the critical value below refer to these same exact objects. Repeated interval occurrences may widen an enclosure but do not select different hosts.

\subsection{Slice, gauge and critical data}\label{sec:quotient}
Use the positive slice
\begin{equation}
s(z)=u+\xi+\iu V\eta,\qquad
z=(\xi,\eta)\in\R^3\times\R^2,\qquad u^T(u+\xi)>0.\label{eq:slice}
\end{equation}
For $h=u^T\Psi\ne0$, multiply $\Psi$ by $\overline h/|h|$ and read the real displacement from $u$ and the $V$-projected imaginary part. Denote this operation by $\mathcal Q$. It gives the same positive gauge as rotation by $-\arg h$, and
\begin{equation}
\Fh_\mu=\mathcal Q\circ F_{\lambda_c+\mu}\circ s,
\qquad H(\xi,\eta)=(\xi,-\eta).\label{eq:quotient}
\end{equation}
The gauge is a coordinate representation, not an extra native update. We work where all prestages and the gauge denominator are nonzero. Common-phase and conjugation equivariance hold on this regular domain; in quotient coordinates, $H\Fh_\mu=\Fh_\mu H$.

Define the inherited real and imaginary linear data by
\begin{align}
M_u&=\Id_3+\eps\diag(k-u^2)+gL,
 &J_R&=M_u-2\eps\diag(u^2),\label{eq:linear}\\
B_0&=V^TM_uV,
 &B(\lambda)&=(1-9\lambda)B_0.\label{eq:compression}
\end{align}
Let $m_1>m_2$ be the eigenvalues of $B_0$. M2 verifies $0<m_2<m_1<1$, $\operatorname{spec}(J_R)\subset(0,0.7)$, and
\begin{equation}
\lambda_c=\frac{1+1/m_1}{9},\qquad \sigma=9m_1>0.
\label{eq:critical}
\end{equation}
The critical multiplier is $-1$ at $\lambda_c$, with derivative $-\sigma$; the other imaginary multiplier is $-m_2/m_1$. A rounded decimal for $\lambda_c$ is never used as a new exact parameter.

For the critical sign and normalization, take
\begin{equation}
e=\frac{(B_{0,12},\,m_1-B_{0,11})^T}
 {\|(B_{0,12},\,m_1-B_{0,11})\|_2},
\qquad e_\perp=(-e_2,e_1)^T.\label{eq:critical-vector}
\end{equation}
The normalizing denominator is certified nonzero. This fixes orientation consistently with the accepted M2 data.

\subsection{What is inherited from M2 and M3}\label{sec:inherited}
In the accepted Taylor convention, let $q_c=(0,e)$, and let $\mathcal B,\mathcal C$ be the second and third derivative tensors of the quotient at criticality. With $J=D\Fh_0(0)$ and the H-even slaving vector $w_R$ in its real block,
\begin{equation}
w_R=(\Id_3-J_R)^{-1}\mathcal B_R(q_c,q_c),\qquad
c_{\rm nf}=\tfrac16 e^T\mathcal C_I(q_c,q_c,q_c)
             +\tfrac12 e^T\mathcal B_I(q_c,(w_R,0)).\label{eq:cnf}
\end{equation}
The amplitude slaving term is included. M2 certifies
\begin{equation}
0.0865248497217<c_{\rm nf}<0.0865248497219,\label{eq:c-bound}
\end{equation}
and a unique local H-exchanged cycle, up to its parity, for all sufficiently small positive $\mu$ \cite{M2,PG}. Its normalized leading amplitude and corrected cycle multiplier satisfy
\begin{equation}
a^2=\frac{\sigma}{c_{\rm nf}}\mu+O(\mu^2),\qquad
\rho_{\rm cycle}=1-4\sigma\mu+O(\mu^2).\label{eq:m2-laws}
\end{equation}
The multiplier is for the second iterate at the emerging cycle, not at the persistent host. The local attraction is modulo common phase. M2 supplies neither an explicit admissible positive-$\mu$ interval nor a numerical neighborhood radius.

At exact $\mu=1/200$, repaired M3 supplies a closed quotient box $U_+$, defined by the rational strings at
\begin{center}\small
\path{evidence/branch/submitted/sources/m3_accepted_support.json}\\
\path{values.exact_receipts.Up}.
\end{center}
Its positive weights are stored at \path{values.exact_receipts.positive_weights}. In the inherited weighted norm, $\Fh_{1/200}^2$ maps $U_+$ strictly into itself and contracts with bound $4789/5000$. It has a unique fixed point $p\in U_+$, and the separate symmetry certificate gives $\Fh_{1/200}(p)=Hp\ne p$ \cite{M3,PG}.

The normalized Fourier entrance $(1,-1/2-\iu\sqrt3/2,-1/2+\iu\sqrt3/2)$, of squared norm three, enters the trap at certified update 301. This is a sufficient entry index, not a first-entry claim. Subsequent full-state convergence is to one common-phase rotation of the two-cycle. Historical \path{Um} is only an outer enclosure of the one-step image; it is not a second trapping set or an exact image. We use neither it nor a new entrance computation for endpoint identification.

\begin{table}[htbp]\centering\small
\begin{tabular}{P{.23\linewidth}P{.67\linewidth}}\toprule
Notation & Meaning in this note\\\midrule
$z=(\xi,\eta)$ & Original five quotient coordinates\\
$x=(X,y,\nu)$ & Five normalized unknowns ($X\in\R^3$), not the state $z$\\
$t=a^2$ & Squared amplitude, not a native update index; updates use $n$\\
$m_1,m_2$; $\bar\mu$ & Compression eigenvalues; a patch's parameter midpoint\\
$\entC,\entS$; $C_{\rm area}$ & Entire scalar functions; channel-area vector\\
$\mathcal Q$; $Q$ & Positive gauge; a matrix for the separate stability norm\\
$q_{\rm root}$; $q_{\rm dyn}$ & Root-equation contraction bound; dynamical derivative bound\\\bottomrule
\end{tabular}
\caption{Notation distinctions needed to compare the sources.}\label{tab:notation}
\end{table}

\section{A removable equation for the nontrivial branch}\label{sec:removable}
The host $z=0$ solves $G(z,\mu)=\Fh_\mu(z)-Hz=0$ for every $\mu$. A root finder for $G$ can therefore return the equilibrium rather than the desired cycle. The accepted formulation fixes the critical projection before continuation.

For $\nu>0$, put
\begin{equation}
t=\frac\mu\nu,\quad a=\sqrt t,\quad
\xi=tX,\quad\eta=a(e+t y e_\perp),\qquad
x=(X_A,X_B,X_C,y,\nu).\label{eq:normalization}
\end{equation}
For $\mu>0$, $e^T\eta=a>0$, excluding the host. Reversing $a$ gives the H partner. At fixed positive $\mu$, the determinant of $x\mapsto z$ in the oriented imaginary basis $(e,e_\perp)$ is
\begin{equation}
\det D_xz=\frac{t^5}{2\nu}>0.\label{eq:chart-determinant}
\end{equation}
Indeed the real $X$ block contributes $t^3$, while the derivative of $(a,aty)$ with respect to $(y,\nu)$ has determinant $a^2t/(2\nu)=t^2/(2\nu)$. Thus this chart is invertible for fixed positive parameter. No claim of invertibility at the collapsed native state $\mu=0$ is made.

With $\Fh_R$ denoting displacement from $u$ and $\Fh_I$ its two imaginary quotient coordinates, the normalized residual is
\begin{align}
R_R(x,\mu)&=\frac{\Fh_R(tX,a(e+t y e_\perp);t\nu)-tX}{t},\label{eq:RR}\\
R_I(x,\mu)&=\frac{\Fh_I(tX,a(e+t y e_\perp);t\nu)
                         +a(e+t y e_\perp)}{a^3}.\label{eq:RI}
\end{align}
The last argument here is the offset $\mu=t\nu$, not $\lambda$. For positive $\mu$, $R=0$ is equivalent to $G=0$. At onset, equations \eqref{eq:RR}--\eqref{eq:RI} are evaluated through the exact continuation described next, never as $0/0$.

\subsection{Prestage, phase and gauge without argument jumps}\label{sec:regular-formulas}
Write the native input as $u+tX+\iu a b$, where $b=V(e+t y e_\perp)$. Its complete prestage is $p+\iu a q$, with
\begin{align}
p_i&=u_i+t\bigl\{(J_RX)_i-\eps[3u_itX_i^2+u_ib_i^2
                                 +t^2X_i^3+tX_ib_i^2]\bigr\},\label{eq:p-expansion}\\
q_i&=(M_ub)_i-\eps(2u_itX_ib_i+t^2X_i^2b_i+tb_i^3).\label{eq:q-expansion}
\end{align}
These are exact cubic expansions. The constant term $p(0)=u$ uses the exact host equation, not a small computed residual. Define
\begin{equation}
P_{ji}=p_jp_i+tq_jq_i,\quad Q_{ji}=q_jp_i-p_jq_i,
\quad r_i=(p_i^2+tq_i^2)^{1/2}.\label{eq:pair-data}
\end{equation}
Cubing $v_j\overline{v_i}=P_{ji}+\iu aQ_{ji}$ gives the harmonic identity
\begin{equation}
\sin(3[\arg v_j-\arg v_i])
=a\frac{3P_{ji}^2Q_{ji}-tQ_{ji}^3}{r_j^3r_i^3}.\label{eq:harmonic}
\end{equation}
The denominators stay positive. The integer harmonic makes this expression independent of principal-argument branch changes. Set
\begin{equation}
d_i=(\lambda_c+t\nu)\sum_{j\ne i}
\frac{3P_{ji}^2Q_{ji}-tQ_{ji}^3}{r_j^3r_i^3},\qquad \delta_i=a d_i.\label{eq:d-kick}
\end{equation}
Introduce the entire functions, by their power series,
\begin{equation}
\entC(w)=\sum_{n\ge0}\frac{(-w)^n}{(2n)!},\qquad
\entS(w)=\sum_{n\ge0}\frac{(-w)^n}{(2n+1)!}.\label{eq:entire}
\end{equation}
For nonnegative $w$ these are $\cos\sqrt w$ and $\sin\sqrt w/\sqrt w$, with the removable value at zero understood. Their definition does not require a square root of a negative auxiliary argument. The post-sync output is $o+\iu a j$, where
\begin{equation}
\begin{aligned}
o_i&=p_i\entC(td_i^2)-tq_id_i\entS(td_i^2),\\
j_i&=p_id_i\entS(td_i^2)+q_i\entC(td_i^2).
\end{aligned}\label{eq:post-sync}
\end{equation}
Let $h_R=u^To$, $h_I=u^Tj$ and $D=(h_R^2+t h_I^2)^{1/2}$. The exact gauged components are
\begin{equation}
\widetilde o_i=\frac{o_i h_R+t j_i h_I}{D},\qquad
\widetilde j_i=\frac{j_i h_R-o_i h_I}{D}.\label{eq:positive-gauge}
\end{equation}
Here $D=|u^TF_\lambda(s(z))|$. The multiplication is precisely $\overline h/|h|$; its $u$-projection is positive. Thus \eqref{eq:positive-gauge} agrees with the inherited quotient and its derivatives on the certified chart.

\subsection{Exact divided differences and explicit tails}\label{sec:divided}
An arithmetic pair $(f_0,f_1(t))$ represents the identity $f(t)=f_0+t f_1(t)$. Its second component is retained as a function, not discarded as a Taylor remainder. For example,
\begin{align}
(fg)_1&=f_1g_0+f_0g_1+t f_1g_1,\label{eq:dd-product}\\
(f/g)_1&=\frac{f_1g_0-f_0g_1}{g_0g(t)},\qquad
(\sqrt f)_1=\frac{f_1}{\sqrt{f_0}+\sqrt{f(t)}}.\label{eq:dd-quotient}
\end{align}
The positive denominator conditions are checked. At $t=0$, the inherited host and critical compression give the exact identities
\begin{equation}
\widetilde o(0)=u,\qquad V^T\widetilde j(0)=-e.\label{eq:cancellations}
\end{equation}
Consequently the residual can be read without subtraction of rounded constants:
\begin{equation}
R_R=\widetilde o_1-X,\qquad
R_I=V^T\widetilde j_1+y e_\perp.\label{eq:regular-residual}
\end{equation}
The subscript 1 denotes the exact divided component. This construction removes both apparent singularities in \eqref{eq:RR}--\eqref{eq:RI}.

For numerical enclosures write $\entC(w)=1+w\entC_1(w)$ and $\entS(w)=1+w\entS_1(w)$. Their coefficients are respectively $(-1)^{n+1}/(2n+2)!$ and $(-1)^{n+1}/(2n+3)!$, for $n\ge0$. Retaining $n=0,\ldots,17$ gives the uniform omitted-tail bounds
\begin{equation}
\begin{array}{c|cc}
 & \text{value tail} & \text{derivative tail}\\\hline
\entC_1 & 2/38! & 36/38!\\
\entS_1 & 2/39! & 36/39!
\end{array}
\qquad(|w|\le1).\label{eq:tails}
\end{equation}
The first omitted denominators are $38!$ and $39!$. Successive absolute-term ratios are below $1/2$, also after differentiation, yielding these geometric-tail estimates. The derivative tail is multiplied by the derivative of its argument in automatic differentiation. The accepted boxes satisfy $\delta_i^2<0.0823$, safely within the sufficient domain. This bound belongs to the present exact family; it is not the unrelated turning-capture short-kick bound.

\begin{lemma}[Removable normalized equation]\label{lem:removable}
The formulas above extend $R$ real-analytically to a neighborhood of $\mu=0$, $\nu>0$ near the accepted onset root. For positive $\mu$ on the certified chart, their zeros are exactly the oriented nontrivial solutions represented by \eqref{eq:normalization} of $\Fh_\mu(z)=Hz$.
\end{lemma}
\begin{proof}
The cubic expressions, entire functions and exact divided identities remove the zero-amplitude quotients. At onset the square-root and gauge denominators are positive; positivity persists locally. Composition and division by those nonzero denominators therefore give an analytic extension, using $t=\mu/\nu$. The cancellations \eqref{eq:cancellations} identify that extension with \eqref{eq:RR}--\eqref{eq:RI} for $t>0$. The state chart is invertible by \eqref{eq:chart-determinant}, and the equation scalings $t$ and $a^3$ are nonzero. An auxiliary extension to negative $t$ near zero establishes regularity only; it is not a computed native family with negative squared amplitude.
\end{proof}

\section{Anchoring the actual M2 germ}\label{sec:onset}
The limit of the real equation at zero is
\begin{equation}
R_R(x,0)=(J_R-\Id_3)X+\tfrac12\mathcal B_R(q_c,q_c).\label{eq:onset-real}
\end{equation}
The imaginary limit retains the noncritical component and the unfolding:
\begin{equation}
R_I(x,0)=(B(\lambda_c)+\Id_2)y e_\perp-\sigma\nu e
+\mathcal B_I(q_c,(X,0))+\tfrac16\mathcal C_I(q_c,q_c,q_c).\label{eq:onset-imag}
\end{equation}
These coefficients are those of the accepted full quotient, including real-amplitude slaving. From \eqref{eq:cnf} the limiting solution obeys
\begin{equation}
X_0=\frac{w_R}{2},\qquad
\nu_0=\frac{c_{\rm nf}}{\sigma}>0,\qquad
y_0=-\frac{e_\perp^TT_0}{\gamma},\label{eq:onset-root}
\end{equation}
where
\begin{equation}
T_0=\mathcal B_I(q_c,(X_0,0))+\tfrac16\mathcal C_I(q_c,q_c,q_c),
\qquad\gamma=1-m_2/m_1>0.\label{eq:T0}
\end{equation}
The exact projection $e^TT_0=c_{\rm nf}$ gives the normalization of $\nu_0$. In the order $(X,y,\nu)$ the complete limiting Jacobian has block form
\begin{equation}
D_xR(x_0,0)=
\begin{pmatrix}
J_R-\Id_3 & 0 & 0\\
\mathcal L & \gamma e_\perp & -\sigma e
\end{pmatrix},\qquad
\mathcal L h=\mathcal B_I(q_c,(h,0)).\label{eq:limiting-jacobian}
\end{equation}
The lower-left coupling is generally nonzero and has not been omitted. The real block is invertible by the inherited spectrum, and the two lower-right columns are independent because $\gamma,\sigma>0$. Hence the full five-dimensional Jacobian is nonsingular.

The accepted onset refinement encloses a root at $\mu=0$ strictly inside the first uniform patch. The saved enclosures for $w_R/2$ and $c_{\rm nf}/\sigma$ lie inside its corresponding coordinate enclosures, consistent with \eqref{eq:onset-root}. For orientation only, the latter ratio is approximately $0.022181967628$; its displayed digits are not a defining constant. The uniform first interval is exactly
\begin{equation}
I_0=[0,1/25600].\label{eq:first-patch}
\end{equation}

The implicit-function argument for the normalized system, together with the uniform validation of Section~\ref{sec:validation}, produces a unique continuous normalized solution on $I_0$, with $\nu>0$. Reconstruction gives
\begin{equation}
z(\mu)\longrightarrow0,\qquad
\eta/a\longrightarrow e,\qquad \xi/a^2\longrightarrow w_R/2
\quad(\mu\downarrow0).\label{eq:germ-limit}
\end{equation}
This identifies the branch's onset direction and positive unfolding. To identify the germ itself, use M2's local uniqueness up to H exchange. Its odd/even expansion can be reparameterized by the exact critical projection $a=e^T\eta$; then its normalized variables have the limits \eqref{eq:onset-root}. The noncritical part of \eqref{eq:onset-imag} fixes $y_0$, as well as $X_0$ and $\nu_0$. The local uniqueness region for $R$ therefore contains that same M2 germ, with the selected positive parity. Both descriptions agree there.

The uniform interval \eqref{eq:first-patch} carries this identification to a positive numerical endpoint. No numerical starting offset is assumed to lie in M2's unspecified sufficiently-small interval. The asymptotic identities identify the germ, while the exact removable-system certificate supplies the finite range; an unspecified $O(\mu^2)$ term is never used as a numerical error bar.

\section{Uniform continuation and common-uniqueness joins}\label{sec:validation}
\subsection{The parameter-uniform contraction argument}\label{sec:lemma}
Let $I$ be a closed parameter interval, $B=c+[-r,r]\subset\R^5$ an exact rational box with $r_i>0$, and $A$ a fixed rational matrix. The normalized unknowns use the weighted norm $\normr{x}=\max_i|x_i|/r_i$. Let $\mathcal J$ enclose $D_xR$ on $B\times I$ and $E=\Id_5-A\mathcal J$.

\begin{lemma}[Uniform root validation and shared uniqueness]\label{lem:uniform}
Suppose $R$ is continuously differentiable on a neighborhood of the box and parameter interval. If
\begin{equation}
q_{\rm root}:=\max_i\sum_j\sup|E_{ij}|\frac{r_j}{r_i}<1\label{eq:root-norm}
\end{equation}
and $T_\mu(x)=x-A R(x,\mu)$ maps $B$ strictly into itself for every $\mu\in I$, then there is a unique zero in $B$ for each parameter. It depends continuously on the parameter, and smoothly wherever $R$ is smooth. Both $A$ and $D_xR$ are nonsingular on the evaluated region. At a fixed shared parameter, a bound of the same form on a convex region containing two already-known roots identifies those roots; a separate self-map on that region is unnecessary.
\end{lemma}
\begin{proof}
The derivative bound makes $T_\mu$ a contraction in the complete closed box, so Banach's argument gives its unique fixed point. Since $\|\Id_5-A D_xR\|_r<1$, the matrix $A D_xR$ is invertible by its convergent Neumann series. Both square factors are therefore nonsingular. A fixed point of $T_\mu$ is a zero of $R$, and local smooth dependence follows from the implicit-function theorem. The contraction estimate also proves continuity: changes in $T_\mu$ at a fixed state bound root displacement after division by $1-q_{\rm root}$. On a shared convex region, applying the mean-value inequality to two zeros gives $\normr{x-y}\le q_{\rm root}\normr{x-y}$, hence equality of the roots. Their existence is already supplied by the adjacent patches.
\end{proof}

To check the self-map uniformly, let $\bar\mu$ be the midpoint of $I$ and let $\Delta=I-\bar\mu$. With $A$ fixed, integrate $A\partial_\mu R(c,\mu)$ to enclose the center defect. The recorded offset is
\begin{equation}
K-c\ \subseteq\ -A R(c,\bar\mu)
-[A\partial_\mu R(c,I)]\Delta+E[-r,r].\label{eq:center-offset}
\end{equation}
Preconditioning precedes multiplication by $\Delta$, preserving valid cancellations. Strict coordinate inequalities $\sup|(K-c)_i|<r_i$ prove the self-map. Approximate inverses select the finite decimal entries of $A$, subsequently treated as exact rationals; they do not establish nonsingularity by themselves.

The arithmetic evaluates derivatives on outward dyadic supersets of the exact rational box. The enclosure and positivity tests are therefore valid on a slightly larger convex region, while the self-map is tested against the original rational radii. Roots lie in the exact box and hence in its evaluated outer box. The derivative contraction also gives uniqueness in the outer evaluation region. This distinction is used in the connector and refinement inclusions.

\subsection{The finite cover and its proof margins}\label{sec:cover}
The accepted certificate uses the exact partition
\begin{equation}
I_i=[i/25600,(i+1)/25600],\qquad i=0,\ldots,127,\label{eq:partition}
\end{equation}
which covers $[0,1/200]$ without gaps. The normalized box radii are
\begin{equation}
r=(0.0005,0.0005,0.0006,0.001,0.0001).\label{eq:radii}
\end{equation}
These are exact rational choices, not estimated uncertainties. Each patch has its own rational center and preconditioner. All 128 pass the strict self-map test and satisfy $q_{\rm root}<0.614$. Positivity bounds keep $\nu$, the slice projection, the prestage real parts and moduli, and the gauge denominator away from zero. Table~\ref{tab:certificate} collects the coarsened accepted bounds. Their full rational values remain in the retained certificate and summary; the manuscript's checks only transcribe those saved results.

% Generated by helpers/make_tables.py from saved accepted data only.
\begin{table}[htbp]\centering\small
\begin{tabular}{P{.39\linewidth}P{.19\linewidth}P{.32\linewidth}}\toprule
Quantity & Bound / count & Source key\\\midrule
Uniform patches / joins & $128\,/\,127$ & \path{E.patches}, \path{E.joins}\\
Patch root norm & $<0.614$ & \path{E.max_patch_q}\\
Common-uniqueness join norm & $<0.583$ & \path{E.max_join_q}\\
Unfolding variable $\nu$ & $>0.02209$ & \path{E.nu_global[0]}\\
Slice projection $u^T(u+\xi)$ & $>8.884$ & \path{E.minimum_slice}\\
Each prestage real part & $>1.504$ & \path{E.minimum_prestage_real}\\
Gauge modulus $|u^T F|$ & $>8.882$ & \path{E.minimum_gauge_modulus}\\
Squared kick $\delta_i^2$ & $<0.0823$ & \path{E.maximum_squared_kick}\\
Each endpoint $U_+$ margin & $>0.00001098$ & \path{E.endpoint_Up_}\newline\path{coordinate_margins}\\
Two-step dynamical norm (strip only) & $<0.9601$ & \path{S.stability.norm_upper}\\
\bottomrule\end{tabular}
\par\vspace{4pt}{\small\raggedright E: \path{EVIDENCE_SUMMARY.json}; S: \path{ENDPOINT_STABILITY_STRIP.json}. Both are under the accepted \path{submitted/data/} directory.\par}
\caption{Certificate summary. Counts are exact. All displayed decimal bounds are exact rational, conservative comparisons with the saved rational values, not floating estimates. The final row applies only to Proposition~\ref{prop:stability}; the first two norms validate roots, not native dynamical rates.}\label{tab:certificate}
\end{table}


At the shared endpoint of $I_i$ and $I_{i+1}$, a convex connector contains both adjacent evaluated outer boxes. Its weighted derivative defect is below 0.583. The existing endpoint roots are in that common uniqueness region, so Lemma~\ref{lem:uniform} identifies them. There are 127 such separate identification checks. Two boxes merely intersecting would not suffice, since their respective roots could lie in disjoint parts of the boxes. The argument uses common uniqueness, not overlap.

Parameter coverage and root identity are thus different verified ingredients. Local smooth graphs for $x(\mu)$ agree at the joins, including their locally unique continuations. Reconstruction yields a continuous family at zero and a smooth family for $\mu>0$. On positive parameters, the coordinate determinant \eqref{eq:chart-determinant} and the nonzero equation scalings transfer regularity to $G$ in the original quotient coordinates. This proves a graph over $\mu$ for the selected family; it does not classify all other roots in state space.

\section{Endpoint identity and the connection theorem}\label{sec:endpoint}
At exact $\mu=1/200$, the accepted refinement uses a smaller normalized box whose entire evaluated enclosure lies inside the last patch's uniqueness region. It selects the same continued root. Reconstruction through \eqref{eq:normalization} produces an interval in the original five quotient coordinates which lies \emph{strictly inside the exact stored} $U_+$. The smallest coordinate margin is greater than $0.00001098$. No reconstruction from rounded center/radius displays defines this trap.

The accepted positive critical parity already lies in $U_+$; no H swap is needed. Because the root satisfies $\Fh(z)=Hz$ and the quotient commutes with H, it is a two-step fixed point in $U_+$. Inherited M3 uniqueness identifies it with $p$. This uses whole-enclosure containment and uniqueness. An intersection, matching amplitude, close nominal coordinates or a small residual would not identify the endpoint. The outer image enclosure \path{Um} is not used.

\begin{theorem}[Connection of the M2 germ to the exact M3 cycle]\label{thm:connection}
For the exact family \eqref{eq:parameters}, with the host, basis and critical data selected by \eqref{eq:host}--\eqref{eq:critical-vector}, the positively oriented M2 alternating germ extends to a graph $z(\mu)$ for every $0<\mu\le1/200$. It is smooth at positive parameters, continuous to the host $z=0$ as $\mu\downarrow0$, and satisfies
\begin{equation}
\Fh_\mu(z(\mu))=Hz(\mu),\qquad
e^T\eta(\mu)>0,\qquad z(1/200)=p.\label{eq:connection-conclusion}
\end{equation}
Here $p$ is the exact repaired M3 point in $U_+$. The family remains in the certified regular chart. Uniqueness is in the stated certified regions, not among all remote state-space solutions. Zero is an onset anchor, not a nontrivial two-cycle. The theorem does not assert attraction throughout the parameter interval.
\end{theorem}
\begin{proof}
Lemma~\ref{lem:removable} and Section~\ref{sec:onset} identify the normalized onset root with the actual M2 germ, with $\nu_0>0$. The first uniform patch includes zero. Lemma~\ref{lem:uniform}, the exact partition \eqref{eq:partition} and the common-uniqueness connectors extend that root without a gap or ambiguity. The certified positive margins permit the original quotient representation throughout. For $\mu>0$, the state coordinate transformation is nonsingular and $e^T\eta=\sqrt{\mu/\nu}>0$. Smoothness and the host limit follow as above. The endpoint refinement and strict inclusion in $U_+$ identify M3 by its inherited uniqueness theorem.
\end{proof}

\begin{corollary}[Primitive full-state two-cycles]\label{cor:full-state}
For every positive parameter in Theorem~\ref{thm:connection}, the quotient cycle and its lifted native-state cycle have primitive period two. Common-phase rotations give a neutral family of the lifted cycles. The phase-invariant observables $C_{\rm area}$, $W$ and $\Gamma$ reverse sign under the H exchange, at their inherited definitions.
\end{corollary}
\begin{proof}
Since $\eta\ne0$, $Hz\ne z$, and H equivariance gives $\Fh_\mu^2(z)=z$. Thus the quotient period is two. If $A=s(z)$, the gauge relation gives $F(A)=e^{\iu\theta}\overline A$. Regular common-phase covariance and conjugation equivariance imply
\begin{equation}
F^2(A)=e^{\iu\theta}\overline{F(A)}=A.\label{eq:full-lift}
\end{equation}
A full fixed point would project to a quotient fixed point, so the full period is primitive two. This is the inherited M2/M3 lift algebra applied along the newly connected family. Phase invariance and conjugation oddness give the stated sign reversal.
\end{proof}

Sign reversal does not establish a nonzero magnitude at every intermediate parameter; no such all-interval claim for $W$ or $\Gamma$ is made. M3's repaired entrance now lands on a cycle identified with the M2 branch, but no entrance theorem at another $\mu$ follows. Each branch member belongs to a separate autonomous map at its fixed parameter. Continuation describes a family of solutions, not an orbit whose parameter changes at each native update.

\section{A separate, limited attraction theorem}\label{sec:stability}
The contraction of $x\mapsto x-A R(x,\mu)$ validates roots of an auxiliary equation. It is not contraction of the native dynamics. The stability calculation instead differentiates the original cubic prestage, phase formula and positive gauge in the original quotient coordinates. At an H-exchanged root, differentiating equivariance gives
\begin{equation}
P_\mu=D\Fh_\mu(Hz)D\Fh_\mu(z)
      =(H D\Fh_\mu(z))^2.\label{eq:two-step}
\end{equation}
Here H also denotes the diagonal linear involution on the tangent coordinates. The identity is evaluated at the exact symmetric cycle, not at an arbitrary root of the host equation.

Let $Q$ be a selected fixed rational basis matrix and $R_Q$ its selected rational approximate inverse. The accepted exact inverse-defect test proves
\begin{equation}
\epsilon_Q=\|\Id_5-R_QQ\|_\infty<1.\label{eq:inverse-defect}
\end{equation}
Hence $Q$ is nonsingular and
\begin{equation}
\|Q^{-1}P_\mu Q\|_\infty
\le\frac{\|R_QP_\mu Q\|_\infty}{1-\epsilon_Q}.
\label{eq:similarity-bound}
\end{equation}
This follows from $Q^{-1}=(R_QQ)^{-1}R_Q$ and the Neumann bound for $(R_QQ)^{-1}$. The matrices are validated as chosen rationals; their numerical origin is not an inverse proof.

\begin{proposition}[Attraction on the explicit endpoint strip]\label{prop:stability}
For the connected cycle of Theorem~\ref{thm:connection}, the actual two-step quotient derivative satisfies
\begin{equation}
\|Q^{-1}P_\mu Q\|_\infty<0.9601,\qquad
\frac{4999}{1000000}\le\mu\le\frac1{200}.\label{eq:stability-strip}
\end{equation}
The cycle is locally attracting in the quotient throughout this strip, and attracting modulo common phase in the full state. This is not a new trapping-neighborhood or uniform-basin theorem.
\end{proposition}
\begin{proof}
The strip's separate normalized root box is contained in the last continuation patch's uniqueness region. Its uniform root validation therefore identifies the same branch. The accepted direct-derivative enclosure and \eqref{eq:inverse-defect}--\eqref{eq:similarity-bound} yield the strict rational comparison with $0.9601$; the bound's approximate display is $0.9600371642733128$. A derivative norm strictly below one at the two-step fixed point gives local quotient attraction by continuity of the derivative. The full common-phase direction is neutral, so the contraction is stated in the quotient. The inherited regular lift transfers the local conclusion to the cycle family modulo common phase.
\end{proof}

\begin{table}[htbp]\centering\small
\begin{tabular}{P{.28\linewidth}P{.32\linewidth}P{.30\linewidth}}\toprule
Result & Parameter scope & What it does not supply\\\midrule
Connected branch identity & $0<\mu\le1/200$ & Attraction over that whole interval\\
Inherited M2 attraction & Sufficiently small $\mu>0$, upper extent unspecified & A numerical endpoint to join to the strip\\
New explicit attraction & $4999/1000000\le\mu\le1/200$ & A uniform basin or new trap\\
Inherited M3 capture & Exact $\mu=1/200$; sufficient entry 301 & Capture at other parameters\\\bottomrule
\end{tabular}
\caption{Connection and attraction have distinct domains. The explicit new strip has width $10^{-6}$; no numerical length is assigned to M2's qualitative neighborhood.}\label{tab:stability-scope}
\end{table}

The 128 broad stability-box estimates in the retained investigation are inconclusive sufficient bounds. They do not prove instability or a bifurcation. One narrow endpoint-strip check establishes Proposition~\ref{prop:stability}; it was not an optimized threshold search. The unknown extent of M2 attraction is not silently joined to this strip. In particular, attraction throughout the full connected interval remains unresolved, even though the branch-identity question is closed.

\section{Evidence, reproducibility and review scope}\label{sec:evidence}
The mathematical investigation was internally accepted on 11 November 2026, without corrections \cite{Review}. This companion manuscript was prepared and completed GPT project-internal manuscript review on 12 November 2026. The author has authorized repository publication of research edition v0.1. Historical submission and acceptance statuses remain unchanged in the preserved evidence.

\subsection{Trusted basis and what was actually checked}\label{sec:trust}
The accepted arithmetic uses outward dyadic intervals at scale $2^{-240}$, exact rational conversion, outward products and reciprocals, integer square-root bounds, exact divided differences and the explicit entire-function tails. Derivatives are forward-mode derivatives of these expressions, with their tail enclosures; floating finite differences do not define any certified derivative. The inherited host is constrained by its accepted box and exact equation. A rounded host midpoint is not substituted into the mathematical premises.

The investigation's generator and checker share the new interval and automatic-differentiation engine. The checker separately organizes exact \texttt{Fraction} comparisons for decisive containment, norm, coverage, join, inverse and endpoint inequalities. This distinction matters: shared-engine replay corroborates execution but is not a second independent interval implementation. The basic dyadic rounding design and the earlier M3 derivative design were consulted as documented references; those historical scripts were not imported or executed for the new investigation.

The accepted original Windows run records 128 validated patches, 127 common-uniqueness joins, the onset and endpoint refinements, four rejected controls and 31 supporting exact algebra/source checks. The controls test false coverage, a broken connector, substitution of the host and an invalid endpoint trap. These are verification counts, not new theorem counts. Nominal planning selected centers and preconditioners in 130 root solves, with 389 Newton iterations in total and at most three at any node. No trajectory was used to bridge the parameter interval.

The subsequent GPT review replayed the submitted checker on Linux Python 3.13.5, changing only the environment guard in a separate replay copy. It did not repeat a Windows execution or implement another interval engine. Five separately formulated high-precision direct-form residual/Jacobian spot checks were numerical diagnostics; they did not certify unsampled parameters. The reviewer artifacts, environment adaptation and unchanged submitted bytes remain in the supplement. No new scientific execution was performed for this manuscript: current Python work only verifies file identities, reads saved data for tables, checks rational transcription and builds or inspects the document.

\subsection{A readable and executable evidence path}\label{sec:source-map}
The complete acceptance extraction is at \path{evidence/branch/}, with the unchanged frozen acceptance ZIP and sidecar at \path{archives/}. The current \path{PRESERVED_INPUTS.json} maps all supplied paths to their delivery paths without changing bytes. The source/claim crosswalk records all six I/N records, actual final equation and theorem labels, pages, and accepted source locations.

For the central evidence, use the following relative entry points; the printed argument above does not require executing them:
\begin{itemize}
\item \href{evidence/branch/submitted/PROOF_REPORT.md}{Accepted proof report}, Sections 2--5: removable formulation, anchor, continuation, joins and endpoint identity.
\item \href{evidence/branch/submitted/data/CONNECTION_CERTIFICATE.json}{Connection certificate}: \path{constants}, \path{onset}, \path{patches}, \path{joins}, \path{endpoint}.
\item \href{evidence/branch/submitted/data/ENDPOINT_STABILITY_STRIP.json}{Narrow stability certificate}: \path{root_certificate}, \path{parent_patch}, \path{stability}.
\item \href{evidence/branch/submitted/data/CERTIFICATE_CHECK.json}{Original checker record} and \href{evidence/branch/submitted/data/ALGEBRA_AND_SOURCE_CHECKS.json}{algebra/source record}; scripts are retained with their historical execution logs.
\item \href{evidence/branch/NATIVE_ALTERNATING_BRANCH_CONNECTION_01_GPT_REVIEW.md}{GPT acceptance review} and \href{evidence/branch/GPT_ACCEPTANCE.json}{machine-readable disposition}; reviewer evidence is under \path{evidence/branch/reviewer/}.
\end{itemize}
Full machine-readable endpoints stay in those records, rather than a printed fraction dump. Saved failed trial bounds and all inconclusive broad stability records are retained. Checksums bind file identities; they do not replace the mathematical argument. The repository checkpoint recorded by the sources is a reading locator, not an instruction to reset a checkout or a fresh claim about its current state.

The new document helpers and build guide are separate from historical scientific scripts. The delivery contains the selectable-text PDF, editable LaTeX, table source and its exact transcription record, reference and font checks, final-page inspection records, source inventories and archive audit. It distributes no font binaries or build caches. Paper G v0.1.1 and all earlier published editions remain unchanged.

\section{Conclusion}\label{sec:conclusion}
The removable onset formulation reaches the actual M2 germ, uniform patches extend it across the full specified interval, common-uniqueness regions identify neighboring roots, and strict endpoint containment identifies the exact M3 cycle. These steps now connect the previously separate local and captured alternating regimes of the same native triad. Inherited symmetry supplies the primitive full-state two-cycles along that family.

The separate attraction result covers only $0.004999\le\mu\le0.005$, in addition to the inherited qualitative near-onset result and endpoint trap at their own scopes. Full-range attraction remains unresolved and is not an additional task concealed in this paper. No global uniqueness or basin, new entrance theorem away from the endpoint, connection to the separate $\lambda=0.5$ orbit, all-interval nonzero observable magnitude, physical field, calibrated time or material-motion law follows.

\textbf{Authorship and review.} Hilmir Fr\'imann Halld\'orsson is the research author. GPT provided scientific steering and project-internal acceptance review; Codex performed the investigation, implementation and manuscript preparation. The mathematical investigation and manuscript have completed project-internal review. The author has authorized repository publication of research edition v0.1. No external peer review, proof-assistant verification or DOI registration is claimed.

\begin{thebibliography}{9}
\bibitem{PG} Hilmir Fr\'imann Halld\'orsson (2026).
\emph{Geometry-Attached Axial Observables, Period-Doubling, and Certified Entrance Capture in the Tri-Octagon Map}.
Paper G v0.1.1, 4 October 2026. Retained \href{evidence/branch/submitted/packet/inputs/paper_g/paper_g.tex}{source text} and \href{evidence/branch/submitted/sources/PAPER_G_README.md}{publication disposition}. Sections 4--6 contain the inherited M2/M3 statements and the earlier non-continuation boundary.

\bibitem{M2} GPT, with Claude and Codex review (2026).
\emph{Trioctagon Magnetism M2: Joint Acceptance Record v0.1}.
2 October 2026. \href{evidence/branch/submitted/sources/M2_JOINT_ACCEPTANCE.md}{Final joint record}, with retained correction and assent records. Internal acceptance of the local theorem; no external peer-review claim.

\bibitem{M3} GPT, with Claude and Codex review (2026).
\emph{Trioctagon Magnetism M3: Joint Acceptance Record v0.1}.
4 October 2026. \href{evidence/branch/submitted/packet/inputs/paper_g/M3_JOINT_ACCEPTANCE.md}{Repaired fixed-parameter acceptance}, with \href{evidence/branch/submitted/sources/m3_accepted_support.json}{exact trap and supporting receipts}.

\bibitem{AB} Hilmir Fr\'imann Halld\'orsson; Codex investigation under GPT scientific steering (2026).
\emph{Native Alternating-Branch Connection: Proof and Verification Report}.
11 November 2026. \href{evidence/branch/submitted/PROOF_REPORT.md}{Accepted investigation report}, original claim ledger, certificate, source inputs and execution records, task \texttt{NATIVE\_ALTERNATING\_BRANCH\_CONNECTION\_01}.

\bibitem{Review} GPT (2026).
\emph{Native Alternating-Branch Connection: Scientific Review and Acceptance}.
11 November 2026. \href{evidence/branch/NATIVE_ALTERNATING_BRANCH_CONNECTION_01_GPT_REVIEW.md}{Internal review}, \href{evidence/branch/GPT_ACCEPTANCE.json}{acceptance disposition} and retained reviewer evidence. No mathematical correction required; narrow attraction scope and shared-engine qualifications retained.

\bibitem{TC} Hilmir Fr\'imann Halld\'orsson (2026).
\emph{Certified Capture of the Collective-Turning Trajectory in the Native Triad}.
Research edition v0.1, 11 November 2026. Retained \href{evidence/branch/submitted/packet/inputs/current/CAPTURE_COMPANION_main.tex}{source} and \href{evidence/branch/submitted/packet/inputs/current/CAPTURE_PUBLICATION_RECEIPT.json}{publication receipt}. Different coefficient family and initial-condition question.
\end{thebibliography}

\end{document}
